% Source: https://artofproblemsolving.com/community/c3h1917978_reflections_about_a_point
% Stable identity: aops-topic-1917978-q01
(a) If $P=(x,y)$ and $P'=(x',y')$, then $M$ is the midpoint exactly when $\dfrac{x+x'}{2}=a$ and $\dfrac{y+y'}{2}=b$, i.e.\ $P'=(2a-x,\ 2b-y)$. The vector from $M$ to $P'$ is $(a-x,\ b-y)$, the negative of the vector from $M$ to $P$. So $P'$ is at the same distance from $M$ as $P$, in the opposite direction: this is exactly the rotation of $P$ through $180^\circ$ about $M$. (No coordinates are needed to construct it: draw the ray from $P$ through $M$ and mark $P'$ beyond $M$ with $MP'=MP$.)

(b) If $P=(x_1,y_1)$ and $Q=(x_2,y_2)$, the vector from $P'$ to $Q'$ is
\[
\bigl((2a-x_2)-(2a-x_1),\ (2b-y_2)-(2b-y_1)\bigr)=-(x_2-x_1,\ y_2-y_1),
\]
the negative of the vector from $P$ to $Q$. So $P'Q'$ is parallel to $PQ$ (pointing the opposite way) and has the same length.

(c) Using $(x,y)\mapsto(8-x,\ 2-y)$:
\[
\boxed{A'(7,1),\quad B'(5,0),\quad C'(7,-1)} .
\]

\begin{asy}
usepackage("amsmath");
size(7cm);
defaultpen(fontsize(10pt));
pair A=(1,1),B=(3,2),C=(1,3),M=(4,1);
pair A2=2*M-A,B2=2*M-B,C2=2*M-C;
draw((-0.5,0)--(8.5,0),gray,Arrow(6)); draw((0,-1.6)--(0,3.8),gray,Arrow(6));
label("$x$",(8.5,0),E,gray); label("$y$",(0,3.8),N,gray);
draw(A--A2,dashed); draw(B--B2,dashed); draw(C--C2,dashed);
filldraw(A--B--C--cycle,paleblue,blue+linewidth(1));
filldraw(A2--B2--C2--cycle,palered,red+linewidth(1));
dot(M); label("$M$",M,SW);
label("$A$",A,SW); label("$B$",B,N); label("$C$",C,NW);
label("$A'$",A2,E); label("$B'$",B2,S); label("$C'$",C2,SE);
label("\mbox{}",(min(currentpicture,user=true)+max(currentpicture,user=true))/2,black); // keeps the text after the figure black when LaTeX typesets the labels inline
\end{asy}

Each dashed segment $AA'$, $BB'$, $CC'$ has midpoint $M$, and the image triangle is the original one turned through $180^\circ$ about $M$.
